\section{Data}
\label{sec:data}

Our dataset comprises 30-minute intraday observations of the S\&P~500
beginning in January 2005. The forecasting target is realized
\emph{variance}, measured as the sum of squared 5-minute log-returns within
each 30-minute bar ($RV_t = \texttt{sumret2}$; see
Equation~\eqref{eq:rv}). We use ``realized variance'' throughout for the
target itself and reserve ``volatility'' for its square root, on which the
models are estimated (Section~\ref{sec:semantic_transforms}).

\subsection{Exogenous Variables}
\label{sec:exog_vars}

In addition to the target series, we incorporate variables spanning several
information categories:

\begin{itemize}
    \item \textbf{Higher-order return moments}: bipower variation, realized semi-variance, autocovariance, absolute returns, cubic returns, and quartic returns.
    \item \textbf{Liquidity and microstructure}: trading volume, turnover (total, buy-initiated, sell-initiated), effective spread, quoted spread, and number of trades.
    \item \textbf{Cross-sectional stock factors}: equal-weighted and value-weighted aggregates of stock-level return moments, bipower variation, semi-variance, and absolute returns.
    \item \textbf{S\&P 500 and sentiment}: SPY trading activity (turnover, buy/sell turnover) and StockTwits social media metrics (attention, sentiment score, sentiment count).
    \item \textbf{Implied volatility}: VIX, VVIX, and VIX3M.
    \item \textbf{Volatility demand}: Options order flow indicators (SPX open/close, all options open/close, open only).
\end{itemize}

\subsection{Feature Accounting}
\label{sec:feature_accounting}

Because several claims in this paper are per-column
(Section~\ref{sec:fwl}) and because the design width is itself part of the
dense-but-weak argument, the mapping from raw series to design columns must be
explicit rather than implied. Three counts circulate in this project---41
exogenous channels, a 359-column exogenous block, and a 529-column design
matrix---and they cannot all be reconciled by the stated expansion rule of
$|\mathcal{L}| = 6$ rolling means per channel, since $41 \times 6 = 246$.
\pending{a per-bucket table giving channels, the lag set applied to each,
derived columns, and the reconciliation to 529; no per-column claim in this
paper should be read as final until this table exists, and the ratio in
Equation~\eqref{eq:per_feature} depends on it directly}

\subsection{Data Cleaning}
\label{sec:cleaning}

The raw data is gridded to a regular 30-minute frequency and filtered to
exclude non-trading periods: Friday evenings after 20:00, all of Saturday, and
Sunday mornings before 18:30. Overnight gaps in stock factors and volatility
demand measures are forward-filled with neutral values. After cleaning and lag
feature generation---which requires a burn-in of 3{,}125 bars for the longest
lag on the base-5 ladder---the frozen evaluation panel comprises exactly
$218{,}934$ observations, and every arm reported in this paper is scored on
those identical rows.

\subsection{A Defect in That Fill, and What It Cost}
\label{sec:fill_defect}

The sentence above---``forward-filled with neutral values''---understates what
the pipeline did, and the understatement was expensive enough to report as a
finding rather than a correction.

The availability indicator that accompanies each imputed exogenous feature was
computed \emph{after} the forward fill. That makes it record ``this series has
started'' rather than ``this bar is observed'', because the fill carries the
last value across every gap. Measured on the volatility-demand family: raw
coverage $0.223$, indicator mean $0.698$. The indicator was zero only for the
leading years before the feed begins.

That single ordering defeated three guards this codebase already contained for
exactly this pathology. The impute-and-indicate median fill never fired,
because the forward fill had left no missing value to fill. The model could not
distinguish a fresh print from a stale one, the indicator being $\approx 1$
either way. And the scale guard---which floors each feature's rolling
interquartile range at an expanding IQR over its \emph{real} values---computed
that floor over the forward-filled constants, so the floor was itself
degenerate and the rolling IQR still collapsed.

The consequence is visible directly in the cached design, whose columns are
expressed in rolling-IQR units and are therefore their own anomaly detector.
The volatility-demand columns reach $2260$, $1547$, $252$ and $206$ IQRs; raw
coverage is $0.223$ and falls to zero through 2024 while the cache continues to
manufacture finite values in the hundreds. Twenty-nine of forty-one channels
exceed $20$ IQRs somewhere and twenty-three pair that with sparse coverage or a
long flat run.

Sparse coverage alone is not the trigger. \texttt{vix}, \texttt{vvix} and
\texttt{vix3m} have the lowest coverage in the panel ($0.159$ to $0.352$) and
the lowest magnitudes ($3.5$ to $8.1$): they are smooth level series whose
filled history retains genuine dispersion. The pathology requires sparse
coverage \emph{and} a collapsing scale---a zero point mass or a flat run---
which is why a signed order-flow series with $15\%$ zeros detonates and an
index level does not.

What it cost is quantified in Section~\ref{sec:era_stability}: on 13--15
October 2023 the volatility-demand contribution swung to $-4.145$ against a
backbone of $+0.928$, driving the forecast one to two orders of magnitude below
the realised value, and $27$ bars carry $47\%$ of the only catastrophic
exogenous failure in twenty years. The fix---record availability before the
fill, bound staleness at $26$ bars, apply the neutral median to unobserved
rows---is in the pipeline; the affected results are marked where they appear.
\pending{re-run of the full battery on the corrected cache, and a restatement
of the \texttt{vol\_demand} bucket effect on bars where the series is actually
observed}

Two properties of the resulting index are worth stating because they bear on
the kernel estimates of Section~\ref{sec:descriptive}. First, observations are
indexed sequentially across the filtered grid, so the long lags span overnight
and weekend boundaries and mix session with non-session bars; the ladder is a
lag structure in bars, not in continuous time. Second, the filter removes
weekend blocks rather than rescaling around them, so the effective sampling of
the kernel is not uniform in calendar time at the longest rungs.
